\cleardoublepage
\phantomsection
\pdfbookmark{Convenzioni tipografiche}{convenzioni-tipografiche}

\begingroup
\let\clearpage\relax
\let\cleardoublepage\relax

\chapter*{Convenzioni tipografiche}

Questa sezione riepiloga le convenzioni tipografiche adottate nel presente lavoro, in particolare per quanto riguarda i riferimenti bibliografici e normativi, l'uso del corsivo e la segnalazione delle parti ancora da completare.

\paragraph{Riferimenti a fonti giuridiche.} Le leggi, i decreti legislativi e le direttive europee citati nel testo (ad esempio la legge Stanca, il D.Lgs.~82/2022 o le direttive (UE) 2019/882 e (UE) 2016/2102) non seguono lo stile di citazione numerico adottato per le altre fonti, ma la convenzione propria della dottrina giuridica italiana: la nota a piè di pagina riporta per esteso \emph{tipo di atto, data completa e numero} (ed eventualmente l'articolo, il comma o l'allegato richiamati), ad esempio:
\begin{quote}
\itshape Legge 9 gennaio 2004, n.~4, art.~3, comma~1.
\end{quote}
Questa scelta rispecchia il modo in cui le fonti normative sono abitualmente richiamate nella letteratura giuridica e nei documenti ufficiali (AGID, Commissione europea), dove l'atto è identificato dagli estremi normativi e non da un numero progressivo di bibliografia.

\paragraph{Riferimenti alle altre fonti (stile IEEE).} Tutte le fonti non giuridiche (standard tecnici come la norma EN~301549, documenti del W3C, linee guida AGID di natura tecnica, letteratura scientifica, rapporti come il WebAIM Million) sono invece citate secondo lo standard IEEE: la nota a piè di pagina riporta, tra parentesi quadre, il numero progressivo con cui la fonte compare in Bibliografia o in Sitografia, seguito dall'eventuale indicazione puntuale (capitolo, paragrafo, pagina), riportata invece fuori dalle parentesi, ad esempio:
\begin{quote}
\itshape [2], cap.~4, § 4.2.
\end{quote}
dove \emph{2} è il numero della fonte in bibliografia. Le fonti bibliografiche in senso stretto (testi, articoli) sono elencate in \emph{Bibliografia}; le pagine web e i documenti consultati online sono elencati separatamente in \emph{Sitografia}, con numerazione progressiva unica e condivisa fra le due sezioni.

\paragraph{Termini stranieri.} I termini e le espressioni in lingua inglese privi di un equivalente italiano consolidato nel linguaggio tecnico (ad esempio \emph{screen reader}, \emph{backend}, \emph{prompt}, \emph{framework}) sono riportati in corsivo in tutte le occorrenze. Non sono invece messi in corsivo gli acronimi (WCAG, LLM, API) né i termini ormai entrati nell'uso corrente italiano.

\paragraph{Termini del glossario.} I termini tecnici che ricevono una definizione puntuale nel Glossario sono contrassegnati, alla loro comparsa nel testo, da un apice \textsuperscript{\textbf{\textit{G}}} (ad esempio \emph{termine}\glox). Il Glossario e l'elenco degli Acronimi e abbreviazioni si trovano in fondo al documento, prima della Bibliografia.

\paragraph{Citazioni testuali.} Le citazioni testuali tratte da fonti normative o tecniche sono racchiuse fra caporali (\enquote{esempio}) per distinguerle dal corsivo usato per i termini stranieri e dall'apice usato per i rimandi al glossario.

\paragraph{Adattamenti all'interno di una citazione testuale.} Quando, all'interno di una citazione testuale, è necessario modificare un elemento grammaticale (ad esempio un cambio di persona, di numero o di tempo verbale) per adattare la citazione alla struttura sintattica della frase che la contiene, la porzione modificata rispetto al testo originale della fonte è racchiusa fra parentesi quadre, secondo la convenzione filologica standard, ad esempio:
\begin{quote}
\itshape \enquote{i componenti dell'interfaccia utente e la navigazione [sono] operabili}
\end{quote}
Le parentesi quadre segnalano quindi un intervento del redattore rispetto al testo originale, senza alterarne il significato.

\paragraph{Codice sorgente.} Gli estratti di codice sono composti in carattere monospazio con evidenziazione sintattica, numerati progressivamente per capitolo ed elencati nel relativo Elenco dei codici sorgenti.

\endgroup